% !TEX root = latex.tex
\documentclass{article}
\usepackage{amsmath}
\usepackage{amsthm}
\usepackage{amssymb}
\usepackage[left=1.5cm, right=1.5cm, top=2.5cm, bottom=2.5cm]{geometry}
\usepackage{fancyhdr} % 引入页眉页脚宏包
\usepackage[hidelinks]{hyperref}
\usepackage{graphicx}
% --- 快捷键设置 ---
\newcommand{\abs}[1]{\left| #1 \right|}

% Mathematical Notation
\newcommand{\R}{\mathbb{R}}
\newcommand{\N}{\mathbb{N}}
\newcommand{\E}{\mathbb{E}}
\newcommand{\mgf}[2]{M_{#1}(#2)}
\newcommand{\Var}{\mathrm{Var}}
\newcommand{\Cov}{\mathrm{Cov}}
\renewcommand{\P}{\mathbb{P}} 
\newcommand{\sig}[2]{\sigma_{#1}^{#2}}
\newcommand{\SD}{\mathrm{SD}}


\newcommand{\sqr}[1]{{#1}^2}
\newcommand{\pow}[2]{{#1}^{#2}}
\newcommand{\nder}[3]{\frac{\mathrm{d}^{#1} #2}{\mathrm{d} #3^{#1}}}
\newcommand{\npder}[3]{\frac{\partial^{#1} #2}{\partial #3^{#1}}}
\newcommand{\uint}[4]{\int_{#1}^{#2} #3 \dd #4}
\newcommand{\dd}{\mathop{}\!\mathrm{d}}
\newcommand{\evalat}[2]{\left. #1 \right|_{#2}}
\newcommand{\eval}[3]{\left. #1 \right|_{#2}^{#3}}
% --- 页眉页脚配置 ---
\pagestyle{fancy}
\fancyhf{} % 清空默认的页眉页脚
\fancyhead[L]{\leftmark} % 左侧显示当前标题
\fancyhead[R]{Saydric Zheng} % 右侧显示你的名字
\fancyfoot[C]{\thepage} % 底部居中显示页码

% 重新定义抓取逻辑：屏蔽 section，只让 subsection 显示在页眉
\renewcommand{\sectionmark}[1]{} 
\renewcommand{\subsectionmark}[1]{\markboth{\thesubsection\ #1}{}}

\begin{document}
\title{Probability and Statistics}
\author{Saydric Zheng}
\date{\today}
\maketitle

\tableofcontents
\newpage
\markboth{}{}


\section{Univariate Discrete Probability Distribution}


\noindent Random Variable: $X$ is a function $(X: S \rightarrow \R)$ that assigns a real number to each point in a sample space $S$.\\

\noindent Discrete Random Variable: a set is countable if its element placed in one-to-one.\\

\noindent Probability Mass Function: $f(x) = \P(X=x)$ for all $x \in \R$. The set of pair $\{(x, f(x)): x \in \R\}$
is the probability distribution \\

\noindent PMF properties:
\begin{itemize}
    \item[(1)] $f(x) \ge 0$ for all $x \in \R$
    \item[(2)] $\sum_x f(x) = 1 $
\end{itemize}

\noindent Cumulative Distribution Function(CDF): $F(x) = \P(X \le x)$

\noindent CDF properties:
\begin{itemize}
    \item[(1)] $0 \le F(X) \le 1$ for all $x \in \R$.
    \item[(2)] $F(X)$ is non decreasing function for all $X \in \R$. 
    \item[(3)] $\lim\limits_{x \to -\infty} F(X+h) = 0$.
    \item[(4)] $\lim\limits_{h \to \infty}F(X+h) = 1$
    \item[(5)] $\lim\limits_{h \to 0^{+}}F(X+h) = F(X)$.
\end{itemize} 

\noindent \textbf{Remark}: $f(a_i) = F(a_i) - F(a_{i-1})$ for $x \ge 2$ \\\\
\noindent Function of Random Variable: If $X$ is a discrete random variable with range $\mathbb{A}$, then $Y = g(X)$ is also discrete random variabel, and its pmf can be determined using the pmf of X.

$$f_{Y}(y) = \P(Y=y) \quad \text{for any value of } y$$
$$g(x) = y, \{x \in \mathbb{A}: g(x) = y\}, \quad \text{define}: f_{X}(x) = \P(X=x)$$
$$f_{Y}(y) = \sum_{{x \in \mathbb{A}: g(x) = y}} f_{X}(x)$$

\noindent Exmaple: $X$ represent the sum of upturned faces o the two thrown dice. Determine the probability distribution \\ $Y = x \bmod 5$.

$$f_{Y}(y) = \P(Y=y)= \P(X=g(y) ^ {-1}) =f_{X}(g(y)^{-1})$$

\subsection{Expetation of Random Varibale}
\noindent Expected Value: \\
$$\E(x) = \sum_{\text{all } x}xf(x) = \sum_{x \in \mathbb{A}} xf(x)$$
\noindent Remark: expected value might can't gained.\\

\noindent Theorem: Suppose $X$ os discrete random variable with range $A$ and $pmf f(x)$. Then the expected value of some real valued function $g(x)$.
$$\E(g(X)) = \sum_{x \in \mathbb{A}} g(x)F(x)$$

\begin{proof}
$$\E(g(X)) = \E(Y) = \sum_{\text{all } y}yf_{Y}(y)$$ where $f_{Y}(y) = \sum_{x \in \mathbb{A}, g(x) = y} f(x)$
Let $D_{y} = \{x \in \mathbb{A}: g(x) = y\}$, $f_{Y}(y) = \sum_{x \in D_{y}} f(x)$
\begin{align*}
   \E(Y) &= \sum_{\text{all } y} \sum_{x \in D_{y}}yf(x) \\
         &= \sum_{\text{all } y} \sum_{x \in D_{y}}g(x)f(x) \\
         &= \sum_{x \in \mathbb{A}}g(x)f(x)
\end{align*}
\end{proof}

\noindent Theorem: For real constant $c_{1}, c_{2}, ... , c_{n}$ and real valued function $g_{1}, g_{2}, ... ,g_{n}$.
$$ \E(\sum_{i = 1}^{n} c_{i}g_{i}(x)) = \sum_{i = 1}^{n}c_{i}\E(g_{i}(x))$$

\begin{proof}
$$h(x) = \sum_{i = 1}^{n} c_{i} g_{i}(x)$$
Substitute into the expected value formula.
\begin{align*}
\E(h(x)) &= \sum_{x \in \mathbb{A}}h(x)f(x) \\
         &= \sum_{x \in \mathbb{A}} (\sum_{i = 1}^{n}c_{i}g_{i}(x))f(x) \\
         &= \sum_{i = 1}^{n}c_{i} (\sum_{x \in \mathbb{A}}g_{i}(x)f(x)) \\
         &= \sum_{i = 1}^{n}c_{i} \E(g_{i}(x))
\end{align*}
\end{proof}
\newpage

\noindent Corollary: For real constant $a$ and $b$.
$$\E(aX+b) = a\E(X) + b$$
\begin{proof}
Let $c_{1} = a, c_{2} = b, g_{1}(x) = x, \textbf{and}, g_{2}(x) = 1$, then by theorem $ \E(cg(X)) = c\E(X)$
\begin{align*}
    \E(aX + b) &= a\E(X) + b\E(1)\\
               &= a\E(X) + b \sum_{\text{all } x \in \mathbb{A}}1 \cdot f(x)\\
               &= a\E(X) + b
\end{align*}
\end{proof}

\noindent Remark: If $g(x)$ is not linear function , then $\E(g(x)) \ne g(\E(X))$

\noindent Variance: The average square distance from variable x to the mean $\mu$.
$$\Var(X) = \E((x - \mu)^2) \Rightarrow \E((x - \mu)^2) = \sum_{x \in \mathbb{A}}(x - \mu)^2f(x)$$ 

\noindent Standard Deviation: $$sigma = SD(x) = \sqrt{\Var(x)} = \sqrt{\E(X - \mu) ^ 2}$$

\noindent Theorem:    
$$\Var(X) = \E(X^2) - \mu ^ 2$$
$$\Var(X) = \E(X(X - 1)) + \E(X) - \mu ^2$$

\begin{proof}
\begin{align*}
\Var(X) &= \E((X - \mu) ^2) \\
        &= \E(X^2 - 2X\mu + \mu^2) \\
        &= \E(X^2) - 2\mu \cdot \E(X) + \mu^2 \\
        &= \E(X^2) - \mu^2
\end{align*}
\end{proof}

\noindent Theroem: 
$$\Var(aX + b) = a^2 \Var(X)$$
$$SD(aX + b) = |a| SD(X)$$
\begin{proof}
$Y = aX + b$, then $\mu_{Y} = a\mu_{X} + b$ \\
\begin{align*}
\Var(Y) &= \E((Y - \mu_{Y})^2) \\ 
        &= \E([aX + b - (a\mu_{X} + b)] ^ 2) \\
        &= \E([aX - a\mu_{X}]^2) \\
        &= \E(a^2(X - \mu_{X})^2) \\
        &= a^2\E((X - \mu)^2) \\
        &= a^2 \Var(X)
\end{align*}
We known that $SD(x) = \sqrt{\Var(X)}$, we square root both side, then:
$$SD(X) = \sqrt{a^2\Var(X)}$$
$$SD(x) = |a|\sqrt{\Var(X)}$$
\end{proof}

\subsection{Moment Generating Function}
Moment Generating Function:
$$M(t) = \E(e^{tx}) = \sum_{x \in \mathbb{A}}e^{tx}f(x)$$
\noindent Example $f(X) = \frac{12}{35(X + 2)}$, for x = -1, 0, 1, 2. Find $mgf$.
$$M(t) = \E(e^{tx}) = \sum_{x \in \mathbb{A}}{e^{tx}f(x)} = \sum_{x \in \mathbb{A}}\frac{12}{25(X + 2)}$$
\begin{align*}
M(t) &= \frac{12}{25}e^{-t} + \frac{6}{25}e^0 + \frac{4}{25}e^t + \frac{3}{25}e^{2t} \\
     &= \frac{6}{25} + \frac{12}{25}e^{-t} + \frac{4}{25}e^t + \frac{3}{25}e^{2t}
\end{align*} 

\noindent Example:\\\\
Theorem: Uniqueness Theorem of $mgf$
Suppose random variable $X, Y$ have $M_{X}(t), M_{Y}(t)$. If $M_{X}(t) = M_{Y}(t)$ for all $t \in (-a, a)$, for $a > 0$. Then $X, Y$ have same probability distribution

\noindent The quantity of $\E(X^n), n = 1, 2, 3,...,$ is called $n^{th}$ moment of probability distribution, or $n^{th}$ moment to $X$.\\\\
\noindent Theorem: If $x$ is a random variable with $mfg M(t)$, then
$$\E(X^n) = M^{(n)}(0)$$

\begin{proof}
\noindent Assume existence $M(t) = \sum_{\text{all } x}e^{tx}f(x)$ for $t \in (-a, a)$
\begin{align*}
M^{(n)}(t) &= \nder{n}{}{t} (\sum_{x \in \mathbb{A}}^{\infty}e^{tx}f(x)) \\ 
         &= \sum_{\text{all } x} \nder{n}{}{t} e^{tx}f(x) \\
         &= \sum_{\text{all } x} (x^n)(e^{tx})f(x)  
\end{align*}
Let $t = 0$, then all the $e^{tx} = 1$, which implies:
\begin{align*}
M^{(n)}(0) &= \left. \nder{n}{M(t)}{t} \right|_{t = 0} \\
           &= \sum_{\text{all } x}(x^n)(e^{tx})f(x) \\
           &= \sum_{\text{all } x \in \mathbb{A}} x^nf(x) \\
           &= \E(x^n)
\end{align*}
\end{proof}

\noindent Example: $f(X) = (\frac{1}{2})^X$, calculate mean and variance of $X$.
Use $\E(X^n) = \nder{n}{M(0)}{t}$.
\begin{align*}
\mu &= \frac{\mathrm{d}}{\mathrm{d}t}(M(0))\\
    &= \left[ \frac{e^t}{2}(1 - \frac{e^t}{2})^{-1} + \frac{e^{2t}}{4}(1- \frac{e^t}{2})^{-2} \right]_{t = 0}\\
    &=\frac{1}{2}(1 - \frac{1}{2})^{-1} + \frac{1}{4}(1 - \frac{1}{2})^{-2}\\
    &= 2
\end{align*}

\begin{align*}
    \E(X ^ 2) &= M^{(2)}(0) \\
              &= \left[ \frac{e^{2t}}{4}{-2}(1- \frac{e^t}{2})^{-3} + \frac{2e^{2t}}{4}(1 - \frac{e^t}{2})^{-2} + \frac{e^t}{2}(1 - \frac{e^t}{2})^{-1} \right]_{t = 0} \\
              &= \frac{1}{4}(1 - \frac{1}{2})^{-3} + \frac{3}{4}(1 - \frac{1}{2})^{-2} + \frac{1}{2}(1 - \frac{1}{2})^{-1} \\
              &= 6
\end{align*}
$$\Var(X) = \E((X - \mu)^2) = \E(X^2) - \E(X) ^2 = 6 - 2^2 = 2$$\\

\noindent Theorem: Let $X$ be random variable with $mgf M_{X}(t)$. If $a$ and $b$ are real constant such that $Y = aX + b$, them $Y$ has $mgf$ given by $M_{Y} = e^{bt}M_{X}(at)$.

\begin{proof}
Since $Y = aX + b$,
\begin{align*}
    M_{Y} &= \E(e^{tY}) \\
          &= \E(e^{t(aX+b)}) \\ 
          &= \E(e^{atX + bt}) \\
          &= \E(e^{atX})\E(e^{bt}) \\
          &= e^{bt}\E(e^{atX}) \\ 
          &= e^{bt}M_{X}({at})
\end{align*}
\end{proof}

\subsection{Special Discrete  Probability Distribution}
\noindent Discrete Uniform Distribution: $X$ takes on values in the set $A = \{a, a + 1, ... , b\}$ such that $a \le b$. Suppose that all values are equal likely.
$$X \sim DU(a, b)$$
\noindent Probability Distribution: $b - a + 1$ values $X$ with equal probability.
$$f(X) = \frac{1}{b - a + 1} \quad \text{for } x = a, a+1, ... , b$$ 

\subsection{Expectation, Covariance, Correlation}
let $U = g(X,Y)$, then $f_U(u) = P(U = u) = \sum_{g(X,Y) = u} f(X, Y)$
$$\E(U) = \sum_{all(X,Y)}g(X,Y),f(X,Y)$$
Extended to n random variables, $X_1, X_2, \dotsc, X_n, U = g(X_1, X_2, \dotsc, X_n)$.
$$\E(U) = \sum_{all(X_1,X_2, \dotsc, X_n)}g(X_1, X_2, \dotsc, X_n)f(X_1, X_2, \dotsc, X_n)$$

\noindent Example: Suppose that $(X_1, X_2, \dotsc, X_n) \sim MN(n, p_1, \dotsc, p_k)$. For i,j $\in {1, 2, \dotsc, k}$ and $i<j$.Determine $\E(X,Y)$
$$f_{ij}(X_i, X_j) = \frac{n! }{x_i!x_j!(n - x_j - x_j)!} (p_i)^{x_i}(p_j)^{x_j}$$
\begin{align*}
    \E(X_iX_j) &= \sum_{all(x_1, \dotsc, x_n)}f(x_ix_j) \\
    &= \sum_{all (x_1, \dotsc, x_n)}x_ix_j(\frac{n!}{x_i!x_j!(n - x_i - x-j)!}(p_i)^{x_i}(p_j)^{x_j}(1-p_i-p_j)^{n - x_i - x_j}) \\
    \intertext{Let $Y_i = X_i - 1$ and $Y_j = X_j - 1$, then:}
    &= n(n-1)p_ip_j\sum_{all(y_i, y_j = 0, 1, \dotsc, n, y_i+y_j \le n - 2)}\frac{(n-2)!}{y_i!y_j!((n-2) - y_i - y_j)} p_i^{y_i}p_j^{y_j}(1 - p_i - p_j)^{((n-2) - y_i - y_j)}\\
    &= n(n-1)p_ip_j(p_i+p_j+(1-p_i-p_j))^{n-2}\\
    &= n(n-1)p_ip_j\\
\end{align*}
For $c_1, c_2, \dotsc, c_n$ are real constants, then
$$\E(\sum_{i=1}^{n}c_ig_i(X_1, X_2, \dotsc, X_n)) = \sum_{i=1}^{n}c_i\E(X_i)$$
\noindent Suppose that $X$ and $Y$ are independent random variables, then
$$\E(g_1(X)g_2(Y)) = \E(g_1(X))\E(g_2(Y))$$
\begin{proof}
    \begin{align*}
            \E(g_1(X)g_2(Y)) &= \sum_{all X,Y}g_1(x)g_2(x)f(x,y)\\
&= \sum_{all X, Y}(g_1(X)f(x))(g_2(Y)f(y))\\
&= \E(g_1(x))\E(g_2(y))
    \end{align*}
\end{proof}

\noindent Corollary:
$$\E(\prod_{i=1}^{n}g_i(X_i)) = \prod_{i=1}^{n}E(g_i(X_i))$$
\begin{proof}
    \begin{align*}
        \E(\prod_{i=1}^{n}g_i(X_i)) &= \sum_{all  X_i}\prod_{i=1}^{n}g_i(X_i)\\
        &= \prod_{i=1}^{n} \sum_{all X_i}g_i(X_i)\\
        &= \prod_{i=1}^{n} \E(g_i(X_i))
    \end{align*}
\end{proof}

\newpage
\noindent Example: $X_1, \dotsc, X_n$ be independent random variables where $X_i \sim Geo(p)$ for $i = 1, 2, \dotsc, k$. Determin the probability distribution of $T = \sum_{i = 1}^{k}X_i$
$$M_T(t) = \frac{p}{1- (1 - p)e^t}t = 1, 2, \dotsc, k, t < \ln(1-p)^{-1}$$
\begin{align*}
    M_T(t) &= \frac{p}{1 - (1-p)e^t}, \quad t=1, 2, \dots, k, \quad t < \ln((1-p)^{-1}) \\
     &= \E(e^{tT})\\
     &= \E(e^{(t(X_1+X_2+ \dotsc X_k))})\\
     &= \E(e^{tX_1})\E(e^{tX_2}) \dotsc \E(e^{tX_k})\\
     &= \prod_{i=1}^{k}M_T(t)(e^{tX_i})\\
     &= (\frac{p}{1-(1-p)e^t})^k
\end{align*}

\noindent Covariance: Random Variables $X$ and $Y$,
$$\Cov(X,Y) = \E(XY) - \mu_X\mu_Y$$
\begin{proof}
    \begin{align*}
        \Cov(X, Y) &= \E((X - \mu_X)(Y - \mu_Y))\\
        &= \E(XY - X\mu_Y - Y\mu_X + \mu_X\mu_Y)\\
        &= \E(XY) - \E(X)\mu_Y - \E(Y)\mu_X + \mu_X\mu_Y\\
        &= \E(XY) - \mu_X\mu_Y - \mu_X\mu_Y + \mu_X\mu_y\\
        &= \E(XY) - \mu_X\mu_Y
    \end{align*}
\end{proof}
\noindent Special case when $X$ and $Y$ are independent random variables, then covariance equation becomes:
$$\Cov(X,Y) = \E(XY) - \mu_X\mu_Y = \E(X)\E(Y) - \mu_X\mu_Y = 0$$


\noindent Example: Suppose that $(X_1, X_2, \dotsc, X_n) \sim MN(N; p_1, \dotsc, p_n), i = 1, \dotsc, k$ and $i < j$.
Determine the $\Cov(X_i, Y_j)$.

$$\E(X_iX_j) = n(n-1)p_ip_j$$
$X_i \sim Bin(n,p_i) \Rightarrow \E(X_i) = np_i$ and $Y_j \sim Bin(n, p_j) \Rightarrow$ $\E(Y_j) = np_j$. 

\begin{align*}
\Cov(X_iX_j) &= \E(X_iX_j) - E(X_i)E(Y_j)\\
&= n(n-1)p_ip_j - n^2p_ip_j\\
&= -np_ip_j
\end{align*}

\noindent Corollary: 
$$\Cov(aX, bY) = ab\Cov(X,Y)$$
\begin{proof}
    \begin{align*}
        \Cov(aX,bY) &= \E((aX)(bY)) - \E(aX)\E(bY) \\
        &= ab\E(XY) - ab\E(X)\E(Y)\\
        &= ab\Cov(X,Y)
    \end{align*}
\end{proof}
\newpage
\noindent Correlation Coefficient $\rho$:
$$Corr = \frac{\Cov(X, Y)}{\sigma_X\sigma_Y}$$
The range of $\rho$ is $-1 \le \rho \le 1$
\begin{proof}
    \begin{align*}
        \intertext{We defined a new variable} Z = Y-tX \intertext{where t is a real number.}
        \Var(Z) &= \E((Z - \mu_Z)^2)\\
                &= \E(((Y - tX) - (\mu_Y - t\mu_X))^2)\\
                &= \E(((Y - \mu_Y) - t(\mu_X - X))^2)\\
                &= \E((Y - \mu_Y)^2 - 2t(X - \mu_X)(Y - \mu_Y) + t^2(\mu_X - X)^2)\\
                &=  t^2\sigma_X^2 - 2\sigma_X\sigma_Y + \sigma_Y^2
    \end{align*}
Since $Var(Z) \ge 0$ for all $t \in \R$, then $t$ at most have one real root.Use discriminant of this quadratic equation must be less or equal to 0.
$$\Cov(X, Y)^2 - 4\sigma_X^2\sigma_Y^2 \le 0$$
$$\abs{\frac{\Cov(X, Y)}{\sigma_X \sigma_Y}} \le 1 \implies -1 \le \rho \le 1$$
\end{proof}

\subsection{Linear Combination of Random Variables}
If $X_{1}, X_{2}, ... , X_{n}$ are $n$ random variables such that $\mu_{i} = \E(X_{i})$ and $\sigma^2_{i} = \Var(X_{i})$ for $i = 1, 2, ... , n$, then we defined V be the linear combination of $c_{i}X_{i}$
$$V = \sum_{i=1}^{n}c_{i}X_{i}$$
where $c_{1}, c_{2}, ... , c_{n}$ are real constants.
$$\E(V) = \E(\sum_{i = 1}^{n}c_{i}X_{i}) = \sum_{i=1}^{n}c_{i}\mu_{i}$$
Consider another sequence of random variables, $Y_{1}, Y_{2}, ... Y_{r}$ and define
$$W = \sum_{j=1}^{r}d_{j}Y_{j}$$
where $d_{1}, d_{2}, ... , d_{r}$ are real constants and $\alpha_{j} = \E(Y_{j})$ for $j = 1, 2, ... , r$.
\begin{align*}
    \Cov(V, W) &= \E((V - \E(V)) (W - \E(W))) \\
               &= \E\left( \left(\sum_{i=1}^{n} c_{i}(X_{i} - \mu_{i})\right) \left(\sum_{j=1}^{r} d_{j}(Y_{j} - \alpha_{j})\right) \right) \\
               &= \sum_{i=1}^{n} \sum_{j=1}^{r} c_{i}d_{j} \E((X_{i} - \mu_{i})(Y_{j} - \alpha_{j})) \\
               &= \sum_{i=1}^{n} \sum_{j=1}^{r} c_{i}d_{j} \Cov(X_{i}Y_{j})
\end{align*}

\noindent Suppose that $n = r$, $X_{i} = X_{j}$ and $c_{i} = c_{j}$. We set $V = W$. Since $\Cov(V,W) = \Var(V)$.
\begin{align*}
    \Var(V) &= \sum_{i=1}^{n}\sum_{j=1}^{n}c_{i}c_{j}\Cov(X_{i}, X_{j}) \\
            &= \sum_{i=j}^{n}c_{i}^2\Var{X_{i}} + \sum_{i=1}^{n}\sum_{j \ne i, j=1}^{n}\Cov(X_{i}, X_{j})\\
            &= \sum_{i=j}^{n}c_{i}^2\sigma_{i}^2 + \sum_{i=1}^{n}\sum_{j=i+1}^{n}c_{i}c_{j}\Cov(x_{i}, x_{j}) + \sum_{i=1}^{n}\sum_{j=1}^{i-1}c_{i}c_{j}\Cov(x_{i}, X_{j}) \\
            &= \sum_{i=j}^{n}c_{i}^2\sigma_{i}^2 + 2\sum_{i=1}^{n}\sum_{j=i+1}^{n}c_{i}c_{j}\Cov(X_{i}, X_{j})
\end{align*}
For $n=2$, then
$$\Var(X_{1} + X_{2}) = \sigma_{1}^2 + \sigma_{2}^2 + 2\Cov(X_{1}, X_{2})$$
$$\Cov(X_{1} + X_{2} + \dotsc X_{n}) = \E(X_{1}X_{2}\dotsc X_{n}) - \E(X_1)\E(X_2), \dotsc, \E(X_n) = 0$$ as $X_1, X_2, \dotsc, X_n$ are independent random variables.Equation can be simplified as
$$\Var(\sum_{i=1}^{n}X_i) = \sum_{i=1}^{n}\sigma_{i}^2$$

\subsection{Conditional Probability Distribution}
Suppose that $X$ and $Y$ are discrete random variables with joint $pmf f(x,y)$.Let $f_X(x)$ be the marginal $pmf$ of $X$ and let $f_Y(y)$ be the marginal $pmf$ of $Y$. The \textbf{conditional probaility mass function(pmf) of $X$ given $Y = y$} is 
$$f_{X|Y}(x|y) = p(X=x|Y=y) = \frac{P(X=x|Y=y)}{P(Y=y)} = \frac{f(x,y)}{f_Y(y)}$$
provided that $f_Y(y) > 0$ Similarly, the \textbf{conditional $pmf$ of $Y$ given $X=x$}is
$$f_{Y|X}(y|x) = P(Y=y|X=x)=\frac{X=x, Y=y} {P(X=x)} = \frac{f(x,y)}{f_X(x)}$$
provided that $f_X(x) > 0$

\noindent for $f_{Y|X}(y|x) \ge 0$ and $f_{X|Y}(x|y) \ge 0$, moreover 
$$\sum_{all x}f_{X|Y}(x|y) = 1 and \sum_{all y}f_{Y|X}(y|x) = 1$$
If $X$ and $Y$ are independent random variables, then $f(x, y) = f_X(x)f_Y(y)$ for all $(x,y)$ implies that\\ $f_{X|Y}(x|y) = f_X(X) \text{and} f_{Y|X}(y|x) = f_Y(y)$
Expended to $n$ random variables: Suppose that  $X_1, X_2, \dotsc X_n$ are discrete random variables with join $pmf f(x_1, x_2, \dotsc, x_n)$ We can define the joint $pmf$ of $(X_i, \dotsc, X_k)$, given $(X_j, \dotsc, X_n)$ by 
$$f_{X_i, \dotsc, X_k|X_j, \dotsc, X_n}(x_i, \dotsc, x_k|x_j, \dotsc, x_n) = P(X_i = x_i, \dotsc, X_k = x_k|X_j=x_j, \dotsc, X_n=x_n) = \frac{f(x_i, \dotsc, f_n)}{f_{X_j, \dotsc, X_n}(x_j, \dotsc, x_n)}$$

\noindent Conditional mean of $X$ given by $Y$:
$$\E(X|Y) = \sum_{all X}x f_{X|Y}(x|y)$$
If $g(X,Y)$ is a function of random variables $X$ and $Y$, then we can define
$$\E(g(X,Y)|Y=y) = \sum_{all x}g(x,y)f_{X|Y}(x|y)$$

\noindent Suppose that $X_1, \dotsc, X_n$ are discrete random variables with joint $pmf f(x_1, \dotsc, x_n)$.
$$\E(\sum_{i=1}^{n}(X_i)) = \sum_{all x_1} \dotsc \sum_{all x_n}(\sum_{i=1}^{n}X_i) f_{12\dotsc(n-1)n}(x_1, x_2, \dotsc, x_{n-1}, x_n)$$

\noindent If $c_1, \dotsc, c_n$ are real constants, and $g_1, \dotsc, g_n$ are arbitrary real valued functions, then
$$\E(\sum_{i=1}^{n}c_ig_(X_i)|Y=y) = \sum_{i=1}^{n}c_i\E(g_i(X_i)|Y=y)$$
\begin{proof}
    \begin{align*}
        \E\left( \sum_{i=1}^{n}c_ig_i(X_i) \,\middle|\, Y=y \right) 
        &= \sum_{\text{all } x_1} \dotsc \sum_{\text{all } x_n} \left[ \left(\sum_{i=1}^{n}c_ig_i(X_i)\right) f_{X_1 \dotsc X_n | Y}(X_1, \dotsc, X_n|Y) \right] \\
        &= \sum_{\text{all } x_1} \dotsc \sum_{\text{all } x_n} \left[ \sum_{i=1}^{n}c_ig_i(X_i) \frac{f(X_1, \dotsc, X_n, Y)}{f_Y(y)} \right] \\
        &= \sum_{i=1}^{n}\frac{c_i}{f_Y(y)} \left[ \sum_{\text{all } x_i} g_i(X_i)f(x_1, \dotsc, x_n, Y) \right] \\
        &= \sum_{i=1}^{n}\frac{c_i}{f_Y(y)} \sum_{\text{all } x_i}g_i(X_i)f_{X_i, Y}(X_i,Y) \\
        &= \sum_{i=1}^{n}c_i \left[ \sum_{x_i} g_i(X_i) f_{X_i|Y}(x_i|y) \right] \\
        &= \sum_{i=1}^{n}c_i\E\left( g_i(X_i) \,\middle|\, Y=y \right)
    \end{align*}
\end{proof}

\noindent Example: $X \sim Poi(\mu)$ and $Y$ given $X=x$, follows  $Bin(x,p)$ distribution. Find conditional $pmf$ of $X$ given $Y = y$.
$$f_{X|Y}(x,y) = \frac{P(X=x, Y=y)}{P(Y=y)}$$
$$P(Y=y|X=x) = \frac{P(X=x, Y=y)}{P(X=x)}$$
Which implies that
$$P(X=x|Y=y) = P(X=x)P(Y=y|X=x)$$
Since $X\sim Poi(\mu)$ and $Y \sim Bin(x, p)$, we get 
$$P(X=x|Y=y) = \frac{e^{-\mu}\mu^x}{x!}\binom{x}{y}p^y(1-p)^{x-y}$$
\begin{align*}
    P(Y=y) &= \sum_{\text{all } x} P(X=x, Y=y) \\
    &= \sum_{x=y}^{\infty} \left( e^{-\mu} \frac{\mu^x}{x!} \binom{x}{y} p^y (1-p)^{x-y} \right) \\
    &= \frac{e^{-\mu} (\mu p)^y}{y!} \sum_{x=y}^{\infty} \frac{(\mu(1-p))^{x-y}}{(x-y)!} \\
    \intertext{Let $z = x - y$, then the summation becomes:}
    &= \frac{e^{-\mu} (\mu p)^y}{y!} \sum_{z=0}^{\infty} \frac{(\mu(1-p))^z}{z!} \\
    &= \frac{e^{-\mu} (\mu p)^y}{y!} e^{\mu(1-p)} \\
    &= \frac{e^{-\mu p} (\mu p)^y}{y!}
\end{align*}
Substitute $P(X=x|Y=y)$ and $P(Y=y)$ into original equation.
\begin{align*}
    f_{X|Y}(x|y) &= \frac{P(X=x, Y=y)}{P(Y=y)} \\
                 &= \frac{\displaystyle \frac{e^{-\mu}\mu^x}{x!} \binom{x}{y} p^y (1-p)^{x-y}}{\displaystyle \frac{e^{-\mu p} (\mu p)^y}{y!}}\\
                 &= \frac{e^{-\mu(1-p)}(\mu(1-p))^{x-y}}{(x-y)!}
\end{align*}

\subsection{Law of Total Expectation}
$$E(g(X)|Y) = E(v(Y)) = \sum_{all y}v(y)f_Y(y) = \sum_{all y}E(g(X)|Y = y)f_Y(y)$$
Theorem: \textbf{Law of Total Expectation}:
Let $g$ be arbitrary real valued function. For random variables $X$ and$Y$. 
$$E(g(X)) = E(E(g(X)|Y))$$
\begin{proof}  
\begin{align*}
    E(E(g(X)|Y)) &= \sum_{all y}E(g(X)|Y=y)f_Y(y)\\
&= \sum_{all y} \left(\sum_{all x} g(X)\frac{f(x,y)}{f_Y(y)}\right) f_Y(y)\\
&= \sum_{all x} g(X)\sum_{all y} f(x, y) \\
&= \sum_{all x} g(X)f_X(x)\\
&=E(g(X))
\end{align*}
\end{proof}

\newpage
\noindent \textbf{Example}: Let $Y$ have the probability distribution:
\begin{center}
    \begin{tabular}{c|ccc}
        $y$ & 1 & 2 & 3 \\ \hline
        $f_Y(y)$ & $1/2$ & $1/4$ & $1/4$
    \end{tabular}
\end{center}
Suppose that $X|Y=y \sim \text{Poi}(\mu_y)$, where $\mu_1=9.4$, $\mu_2=11.9$, and $\mu_3=13.5$. Find $\E(X)$ and $\Var(X)$.
\begin{proof}
By the Law of Iterated Expectations:
\begin{align*}
    \E(X) &= \sum_{y=1}^{3} \E(X|Y=y) f_Y(y) \\
          &= 9.4\left(\frac{1}{2}\right) + 11.9\left(\frac{1}{4}\right) + 13.5\left(\frac{1}{4}\right) \\
          &= 11.05
\end{align*}

\noindent To find $\Var(X)$, we first calculate $\E(X^2)$ using the Law of Iterated Expectations:
\begin{align*}
    \E(X^2) &= \sum_{y=1}^{3} \E(X^2|Y=y) f_Y(y) \\
            &= \E(W_1^2)\left(\frac{1}{2}\right) + \E(W_2^2)\left(\frac{1}{4}\right) + \E(W_3^2)\left(\frac{1}{4}\right)
\end{align*}
Since $\Var(W) = \E(W^2) - (\E(W))^2$, we have $\E(W^2) = \Var(W) + (\E(W))^2$:
\begin{align*}
    \E(X^2) &= \left(\Var(W_1) + \mu_1^2\right)\left(\frac{1}{2}\right) + \left(\Var(W_2) + \mu_2^2\right)\left(\frac{1}{4}\right) + \left(\Var(W_3) + \mu_3^2\right)\left(\frac{1}{4}\right) \\
            &= (9.4 + 9.4^2)\left(\frac{1}{2}\right) + (11.9 + 11.9^2)\left(\frac{1}{4}\right) + (13.5 + 13.5^2)\left(\frac{1}{4}\right) \\
            &= 97.76\left(\frac{1}{2}\right) + 153.51\left(\frac{1}{4}\right) + 195.75\left(\frac{1}{4}\right) \\
            &= 136.195
\end{align*}

\noindent Finally, calculating the variance:
\begin{align*}
    \Var(X) &= \E(X^2) - (\E(X))^2 \\
            &= 136.195 - (11.05)^2 \\
            &= 14.0925
\end{align*}
\end{proof}

\subsection{Expectation, Covariance and Correlation}

$$ \P(U = u) = \sum_{\substack{\text{all }(x, y): \\ g(x, y) = u}} f(x, y) $$
$$ \E(U) = \E(g(X, Y)) = \sum_{\text{all }(x,y)} g(x,y) f(x,y) $$

$$ \bar{x} = \frac{1}{n} \sum_{i=1}^{n} x_i $$
$$ s^2 = \frac{1}{n - 1} \sum_{i=1}^{n} (x_i - \bar{x}) ^ 2$$





\newpage
\section{Univariate Continuous Probability Distribution}
Continuous Random Variable: $P(X=x) = 0$
The properties of a $cdf$ are the same for continuous random variables as for discrete random variables. These properties include:
\begin{itemize}
    \item $F(x)$ is defined for all $x \in \R$,
    \item $F(x)$ is a non-decreasing function of $x$ for all $x\in \R$,
    \item $\lim_{x -> -\infty} F(x) = 0$ and $lim_{x -> \infty} F(x) = 1$
    \item $P(a < x \le b) = F(b) - F(a)$
\end{itemize}

\begin{proof}
    Using property (4), for any $a \in \R$
    $$P(X = a) = \lim_{\epsilon -> 0}P(a-\epsilon < X \le a) = \lim_{\epsilon -> 0}(F(a) - F(a - \epsilon)) = F(a) - \lim_{\epsilon -> 0}F(a - \epsilon) = 0$$
\end{proof}
\noindent Suppose that $X$ is continuous random variable withh $cdfF(x)$. For $p \in (0, 1)$, the $100p^\text{th}$ \textbf{percentile} of $X$ is the value $q(o)$ such that $F(q(p)) = p$. If $p = 0,5$, then $m = q(0.5)$ is called the \textbf{median} of $X$

\subsection{Functions of Random Variables}
\textbf{Example}: Let$X$ be a continuous random variable with $pdf f_{X}(x)$ and $cdf F_{X}(x)$. Determine a general formula for the $pdf$ of $Y=X^2$
Since $Y = X^2$, then $Y$ is non-negative continuous function and $y > 0$.
$$F_{Y}(y) = P(Y \le y) = P(X^2 \le y) = P(- \sqrt{y} \le X \le \sqrt{y}) = F_X(\sqrt{y}) - F_X(- \sqrt{y})$$
\begin{align*}
    f_Y(y) &= F'_Y(y)\\
        &= F'_X(\sqrt{y})\frac{d}{dy}(y^{\frac{1}{2}}) - F'_X(- \sqrt{y})\frac{d}{dy}(-y^{\frac{1}{2}})\\
        &= f_X(\sqrt{y} \frac{1}{2\sqrt{y}}) - f_X(-\sqrt{y})(\frac{-1}{2\sqrt{y}}), \text{for } y > 0\\
\end{align*}

\noindent Theorem: Let $X$ be continuous random variable with $pdf f_X(x)$.Suppose that $g$ is a strictly monotonic differentiable function an range of X. Then, $pdf Y = g(X)$ in range where $f_Y(y) > 0$ given by 
$$f_Y(y) = f_X(g^{-1}(y))|\frac{dg^{-1}(y)}{dy}|$$

\begin{proof}
    Assume that $g$ is monotonically increasing, then
    $$F_Y(y) = P(Y \le y) = P(g(X) \ le y) = P(X \le g^{-1}(y)) = F_X(g^{-1}(y))$$
    $$f_Y(y) = F_Y(y) = \frac{d}{dy}(F_X(g^{-1}(y)) = F'_X(g^{-1}(y)))\frac{dg^{-1}(y)}{dy} = f_X(g^{-1}(y))\frac{dg^{-1}(y)}{dy}$$
    Since $g$ is monotonically increasing, $g^{-1}$ is also monotonically increasing.
    $$|\frac{dg^{-1}(y)}{dy}| = \frac{dg^{-1}(y)}{dy}$$

    \noindent Assume that $g$ is monotonically decreasing,  Then
    $$F_Y(y) = P(Y \le y) = P(g(X) \le y) = P(X \ge g^{-1}(y)) = 1 - F_X(g^{-1}(y))$$
    $$f_Y(y) = F'_Y(y) = \frac{d}{dy}[1 - F_X(g^{-1}(y))] = F'_X(g^{-1}(y))\frac{dg^{-1}(y)}{dy} = f_X(g^{-1}(y))\frac{dg^{_1}(y)}{dy}$$
    Since $g$ is monotonically decreasing, $g^{-1}$ is also monotonically decreasing that implies that the derivatives are negative.
    $$- \frac{dg^{-1}(y)}{dy} = |\frac{dg^{-1}(y)}{dy}|$$
\end{proof}

Corollary: Let $X$ be a continuous random variable with $pdf f_X(x)$. If $a$ and $b$ are real constants such that $a \ne 0$, then $Y = aX + b$ has $pdf$ given by
$$f_Y(y) = \frac{1}{|a|}f_X(\frac{y - b}{a})$$
\begin{proof}
    Define that $y = g(x) = ax + b$, then $x = \frac{y - b}{a}$, then by previous theorem,
    $$f_Y(y) = f_X(g^{-1}(y))|\frac{dg^{-1}(y)}{dy}| = f_X(\frac{y-b}{a}|\frac{d}{dy}(\frac{y-b}{a})|) = \frac{1}{|a|}f_X(\frac{y-b}{a})$$
\end{proof}

\subsection{Expectation of Random Variable}
The expected value or mean of continuous random variable $X$ with $pdf f_X(x)$ is defined by 
$$\mu_X = \E(X) = \uint{-\infty}{\infty}{xf_X(x)}{x}$$

\begin{itemize}
\item $mgf$ of $X$ is defined as:
$$M_X(t) = \E(g(X)) = \uint{-\infty}{\infty}{g(x)f(x)}{x}$$
provided there exists all values of t within some open interval.

\item The $n^{th}$ moment of $X$, where $n = 1, 2, \dotsc$, is defined as:
$$E(X^n) = \uint{-\infty}{\infty}{x^nf_X(x)}{x}$$
$$\E(X^n) = \mgf{X}^{(n)}(0)=\evalat{\nder{n}{\mgf{X}{t}}{t}}{t=0}, \text{for } n = 1, 2, 3, \dotsc$$

\item Variance of $X$ id defined as:
$$\sig{X}{2} = \Var(X) = \E((x-\mu_x)^2) = \uint{-\infty}{\infty}{(x-\mu_k)^2}{x} = \E(X^2) - \E(X)^2$$
$$\sig{X}{} = \SD(X) = \sqrt{\Var(X)}$$

\item For real constant $a$ and $b$.
$$\E(aX+b) = a\mu_X + b$$
$$\Var(aX+b) = a^2\sig{X}{2}$$
$$\SD(aX+b) = \abs{a}\sig{X}{}$$
$$\mgf{aX+b}{t} = e^{bt}\mgf{aX}{t}$$
\end{itemize}



\subsection{Special Continuous Probability Distributions}
\subsubsection{Continuous Uniform Distribution}
Consider a random variable $X$ which takes on values in interval $(a,b)$ where $a$ and $b$ are real numbers such that $a < b$. If the probability distribution of $X$ satisfies all subintervals of $(a,b)$ with the same length are equally likely,
then $X$ is said to have a \textbf{continuous uniform distribution}. Denote as $X \sim U(a, b)$
$$ \uint{a}{b}{f(x)}{x} = \uint{a}{b}{k}{x} = k(b-a) \Rightarrow k = \frac{1}{b-a}$$
$$ F(X) = \uint{a}{\infty}{\frac{1}{b-a}}{x} = \frac{x-a}{b-a}$$
from equation $F(q(p)) = p$
$$\frac{q(p) - a}{b-a} = p \Rightarrow q(p) - a = p(b-a) \Rightarrow = q(p) = a(p-1) + bp$$
$$\E(X) = \uint{-\infty}{\infty}{x(fx)}{x} = \eval{\frac{1}{b-a}\frac{x^2}{2}}{b}{a}
= \frac{\sqr{b} - \sqr{a}}{2(b -a)} = \frac{(b-a)(b+a)}{2(b-a)} = \frac{b+a}{2}$$
\begin{align*}
    \Var(X) &= \E(\sqr{X}) - \sqr{E(X)} \\
    &= \uint{a}{b}{x^2 \frac{1}{b-a}}{x} - \sqr{\frac{b+a}{2}} \\
    &= \frac{1}{b-a} \eval{\frac{x^3}{3}}{a}{b} - \frac{\sqr{b} - 2ab - \sqr{a}}{4} \\
    &= \frac{1}{b-a} \frac{(b-a)(\sqr{b} +ab + \sqr{a})}{3} - \frac{\sqr{b} + 2ab + \sqr{a}}{4}\\
    &= \frac{\sqr{b} - 2ab + \sqr{a}}{12} \\
    &= \frac{\sqr{b-a}}{12}
\end{align*}
Theorem:Suppose that X is continuous random variable with $cdf F_X(X)$. Suppose that $F_X(X)$ is strictly increasing in following case: If $x_1<x_2$ such that $0<F_X(x_1)$ and $F_X(x_2) < 1$, then $F_X(x_1) < F_X(x_2)$ If $Y = F_X(X)$, then $Y \sim U(0,1)$
\begin{proof}
If 
\end{proof}
\subsubsection{Exponential Distribution}
Denote as $X \sim Exp(\theta)$
\[f(x) = \begin{cases}
    \frac{1}{\theta}e^{\frac{-X}{\theta}} & \text{for} x > 0, \\
    0 & \text{otherwise} 
\end{cases}\]
$$\uint{-\infty}{\infty}{f(x)}{x} = \uint{-\infty}{\infty}{\frac{1}{\theta}e^{\frac{-x}{\theta}}}{x} = \eval{-e^{\frac{-x}{\theta}}}{0}{\infty} = \lim_{x -> \infty}-e^{\frac{-x}{\theta}} + 1 = 1$$
$$F(X) = \uint{0}{x}{\frac{1}{\theta}e^{\frac{-z}{\theta}}}{z} = \eval{-e^{\frac{-z}{\theta}}}{0}{x} = 1 - e^{\frac{-x}{\theta}}$$
\[F(X) = \begin{cases}
    0 & \text{otherwise} \\
    1 - e^{\frac{-x}{\theta}} & \text{for} x > 0
\end{cases}\]
Determine the p-quantile, from the equation $F(q(p)) = p:$
$$1 - e^{\frac{-q(p)}{\theta}} = p \Rightarrow \frac{-q(p)}{\theta} = -\ln(p-1) \Rightarrow q(p) = -\theta\ln(1-p)$$

\begin{align*}
    \mgf{}{t} = \E(e^{tx}) &= \uint{0}{\infty}{e^{tx}\frac{1}{\theta}e^{\frac{-x}{\theta}}}{x} \\
        &= \frac{1}{\theta}\uint{0}{\infty}{e^{x(t-\frac{1}{\theta})}}{x} \\
        &= \frac{1}{\theta(t - \frac{1}{\theta})} \eval{\frac{1}{t - \frac{1}{\theta}}e^{x(t- \frac{1}{\theta})}}{0}{\infty}\\
        &= \frac{1}{t\theta - 1} [\lim_{x-> \infty} e^{x(t - \frac{1}{\theta})} - 1]\\
        &= \frac{1}{t\theta - 1} [0 - 1] \indent \text{provided } t -\frac{1}{\theta} < 0 \Rightarrow t < \frac{1}{\theta} \\
        &= \frac{1}{1 - t\theta}
\end{align*}

\noindent $\mgf{1}{t} = \theta(1- t\theta)^{-2}$\\
$\mgf{2}{t} = -2\theta(1-t\theta)^{-3}{-\theta} = 2\sqr{\theta}(1-t\theta)^{-3}$

$$\E(X) = \mgf{1}{0} = \theta$$
$$\Var(X) =\E(\sqr{X}) - \sqr{E(X)} = \mgf{2}{0} - \sqr{\theta} = 2\sqr{\theta} - \sqr{\theta} = \sqr{\theta}$$

\subsubsection{Gaussian Distribution}
$$f(X) = \frac{1}{ \sqrt{2\pi}\sigma} e^{\frac{-1}{2}(\frac{x-\mu}{\sigma})^2} \text{for} -\infty < x < \infty$$
where $\mu \in \R$ and $\sigma > 0$ are parameter of probability distribution.
To verify the area under the curve of $f(X)$ equal 1.
\begin{align*}
    \uint{-\infty}{\infty}{\frac{1}{\sqrt{2\pi}\sigma}e^{\frac{-1}{2}\sqr{(\frac{x-\mu}{\sigma})}}}{x} &= \uint{-\infty}{\infty}{\frac{1}{\sqrt{2\pi}}e^{\frac{-\sqr{y}}{2}}}{y} \indent \text{if } y = \frac{x - \mu}{\sigma} \Rightarrow dy = \frac{dx}{\sigma} \\
    &= \frac{1}{\sqrt{2\pi}} \cdot 2 \uint{0}{\infty}{e^{\frac{-y^2}{2}}}{y} \indent \text{Since } e^{\frac{-y^2}{2}} \text{is an even function}\\
    &= \frac{\sqrt{2}}{\pi}\uint{0}{\infty}{e^{\frac{-\sqr{y}}{2}}}{y}
\end{align*}
\begin{align*}
    \sqr{\left(\frac{\sqrt{2}}{\sqrt{\pi}}\uint{0}{\infty}{e^{\frac{-\sqr{y}}{2}}}{y}\right)} &= \frac{2}{\pi}\left(\uint{0}{\infty}{e^{\frac{-\sqr{y}}{2}}}{y} \uint{0}{\infty}{e^{\frac{-\sqr{z}}{2}}}{z}\right) \\
    &= \frac{2}{\pi}\left(\uint{0}{\infty}{\uint{0}{\infty}{e^{-(\sqr{y} + \sqr{z})}}{z}}{y}\right)\\
    &= \frac{2}{\pi} \uint{0}{\infty}{re^{\frac{-\sqr{r}}{2}}}{r} \indent \text{Let } y = r\cos\theta \text{and } z = r\sin\theta \Rightarrow dydz = rdrd\theta \\
    &= \uint{0}{\infty}{e^{-s}}{s}  \indent \text{Let} s = \frac{\sqr{r}}{2}
    &= \eval{-e^{-s}}{0}{\infty}
    &= 1
\end{align*}

Standard Nornal Distribution:
$$f_Z(Z) = \frac{1}{\sqrt{2\pi}}e^{\frac{-\sqr{z}}{2}} \indent \text{for} -\infty < z <\infty, Z \sim N(0, 1)$$
$mgf$:
\begin{align*}
    \mgf{Z}{t} &= \uint{-\infty}{\infty}{e^{tz}\frac{1}{\sqrt{2\pi}}e^{\frac{-\sqr{z}}{2}}}{z} \\
    &= \uint{-\infty}{\infty}{\frac{1}{\sqrt{2\pi}}e^{\frac{-1}{2}(z^2 - 2tz)}}{z} \\
    &= \uint{-\infty}{\infty}{\frac{1}{\sqrt{2\pi}}e^{\frac{-1}{2}\sqr{(z - t)}} e^{t^2}}{z} \indent \text{Since } \uint{-\infty}{\infty}{\frac{1}{\sqrt{2\pi}}e^{\frac{-1}{2}\sqr{(z-t)}}}{z} \text{is } Z \sim N(t, 1)\\
    &= e^{\sqr{t}} \indent \text{for } t \in \R
\end{align*}
$$\E(Z) = \evalat{\frac{d}{dz}(e^{\frac{\sqr{t}}{2}})}{z=0} = te^{\frac{\sqr{t}}{2}} = 0$$
\begin{align*}
    \Var(Z) &= \E(\sqr{Z}) - \sqr{\E(Z)} \\
    &= \evalat{\nder{2}{}{z}\left[e^{\frac{\sqr{t}}{2}} + \sqr{t}e^{\frac{\sqr{t}}{2}}\right]}{t=0}\\
    &= 1
\end{align*}

Theorem: If $X\sim N(\mu, \sqr{\sigma})$ and $z = \frac{x - \mu}{\sigma}$, then $z \sim N(0,1)$.
\begin{proof}
Express $z = ax + b$ for $a = \frac{1}{\sigma}$ and $b = \frac{-\mu}{\sigma}$
By $$f_Y(y) = \frac{1}{\abs{a}}f_X(\frac{y-b}{a})$$
\begin{align*}
    f_Z(z) &= \frac{1}{\abs{\frac{1}{\sigma}}}f_X(\frac{z + \frac{\mu}{\sigma}}{\frac{1}{\sigma}}) \\
    &= \sigma f_X(z\sigma + \mu) \indent \text{Since } X \sim N(\mu, \sqr{\sigma})\\
    &= \sigma \frac{1}{\sqrt{2\pi \sigma}}e^{\frac{-1}{2} \sqr{(\frac{\sigma z + \mu - \mu}{\sigma})}} \\
    &= \frac{1}{\sqrt{2\pi}} e^{\frac{-1}{2}\sqr{z}} \indent \text{for} -\infty < z < \infty, \text{ such that} z \sim N(0, 1)
\end{align*}
\end{proof}
From previous theorem, $z= \frac{x - \mu}{\sigma}$, $z \sim N(0,1)$, if $X \sim N(\mu, \sqr{\sigma})$
$$\mgf{X}{t} = \mgf{\sigma Z + \mu}{t} = e^{\mu t}\mgf{Z}{\sigma t} = e^{ut} \cdot e^{\frac{\sqrt{(\sigma t)}}{2}} \indent \text{for} t \in \R$$
$$\E(X) = E(\sigma Z + \mu) = \sigma \E(Z) + \mu = \sigma \cdot 0 + \mu = \mu$$
$$\Var(X) = \Var(\sigma Z + \mu) = \sqr{\sigma}\Var(Z) + 0 = \sqr{\sigma} \cdot 1 = \sqr{\sigma}$$

\subsubsection{Central Limit Theorem(CLT)}
Theorem: (Central Limit Theorem(CLT)):Let$X_1, \dotsc, X_n$ be independent random variables that all have same probability distribution, with $\mu$ and $\sqr{\sigma}$.Let $S_n = \sum_{i=1}^{n}X_i$. Then as $n -> \infty$, the $cdf$ of the random variable
$$\frac{S_n - n\mu}{\sigma \sqrt{n}}$$
approaches the $N(0, 1) cdf$.

\noindent Sample Mean: $$\bar{X_n} = \frac{1}{n}\sum_{i=1}^{n}X_i = \frac{S_n}{n}, \text{substitute into equation} \frac{S_n - n\mu}{\sigma\sqrt{n}} $$\\
$$\frac{S_n - n\mu}{\sigma \sqrt{n}} = \frac{n \bar{X_n} - n \mu}{\sigma \sqrt{n}} = \frac{\bar{X_n} - \mu}{\frac{\sigma}{\sqrt{n}}}$$
Remark: Apply CLT when $n$ is large and finite to approximate the proability distribution of $\sum_{i=1}^{n}X_i$ or $\bar{X}$ by normal distribution.
$$S_n = \sum_{i=1}^{n}X_i \sim N(n\mu, n\sqr{\sigma}) \text{ and } \bar{X_n} \sim N(\mu, \frac{\sqr{\sigma}}{n})$$







\section{Introduction to Statistical Science}
Variate: A varaite is a characteristic of a unit.
Attribute: An attribute of a population or process is a function of the variate over the population or process.
Three types of empirical studies:
\begin{itemize}
    \item Sample Survey: Obtained information by selecting a "representative" sample of units form the population and determining the variates of interest for each unit in the sample.
    \item Observational Studies: An observational study is one in which data are collected about a population or process without any attempt to change the value of one or more variates for the sampled units. 
    \item Experimental Studies: An experimental study is one in which the experimenter (that is, the person conducting the study) intervenes and changes or sets the values of one or more variates for the units in the sample. 
\end{itemize}

\noindent Data Summaries:\\
\indent Sample Mean: $\bar{y = \frac{1}{n}\sum_{i=1}^{n}X_i}$\\
\indent Sample Median: $\hat{m}$: the meddle value when n is odd and the sample is ordered from smallest to largest and the average of the two middle values when n is even. \\
\indent Sample Mode: the value of y which have highest frequency appears in the sample. \\
\noindent Measure of dispersion or variability\\
\indent Sample Variance: $$\sqr{s} = \frac{1}{n-1}\left[\sum_{i=1}^{n}\sqr{y_i} - \frac{1}{n}\sqr{\left(\sum_{i=1}^{n}y_i\right)}\right]$$ \\
\indent And the sample standard deviation: $s = \sqrt{\sqr{s}}$\\\\
\begin{proof}
    \begin{align*}
        \sqr{s} = \frac{1}{n-1}\sum_{i=1}^{n}\sqr{\left(y_i - \bar{y}\right)} &= \frac{1}{n-1}\left[\sum_{i=1}^{n}\sqr{y_i} - 2y_i\bar{y} -  \sqr{\bar{y}}\right]\\
            &=\frac{1}{n-1}\left[\left(\sum_{i=1}^{n}\sqr{y_i}\right) - 2\bar{y}\left(\sum_{i=1}^{n}y_i\right) + \sum_{i=1}^{n}\sqr{\bar{y}}\right] \indent \text{Since } \bar{y} = \frac{1}{n}\sum_{i=1}^{n}y_i\\
            &=\frac{1}{n-1}\left[\left(\sum_{i=1}^{n}\sqr{y_i}\right) - 2n\sqr{\bar{y}} - n\sqr{\bar{y}}\right]\\
            &=\frac{1}{n-1}\left[\left(\sum_{i=1}^{n}\sqr{y_i}\right) - n\left(\sqr{\sum_{i=1}^{n}y_i}\right)\right]\\
            &=\frac{1}{n-1}\left[\sum_{i=1}^{n}\sqr{y_i} - \frac{1}{n}\sqr{\sum_{i=1}^{n}y_i}\right]
    \end{align*}
\end{proof}

\noindent Expectation of Sample Mean:
$$\E(\bar{X}) = \E(X)$$

\begin{proof}
    $$\E(\bar{X}) = \E(\frac{1}{n}\sum_{i=1}^{n}X_i) = \frac{1}{n}\sum_{i=1}^{n}\left(\E(X_i)\right) = \frac{1}{n} n \E(X) = \E(X)$$ 
\end{proof}
\noindent Variance of Sample Mean:
$$\Var(\bar{X}) = \frac{\Var(X)}{n}$$
\begin{proof}
    \begin{align*}
        \Var(\bar{X}) &= \Var(\frac{1}{n}\sum_{i=1}^{n}X_i)\\
        &= \frac{1}{\sqr{n}}\Var(\sum_{i=1}^{n}X_i)\\
        &= \frac{\sum_{i=1}^{n}\Var(X_i)}{\sqr{n}}\\
        &= \frac{n\Var(X)}{\sqr{n}}\\
        &= \frac{\Var(X)}{n}
    \end{align*}
\end{proof}
\noindent The stnadard error is $\sigma_{\bar{X}} = \frac{\sigma}{\sqrt{n}}$\\
\noindent Expectation of Sample Variance:
$$\E(\sqr{s_n}) = \Var(X)$$

\begin{proof}
    \begin{align*}
        \E(S_n^2) &= \E\left(\frac{1}{n-1}\sum_{i=1}^{n}\sqr{(X_i - \bar{X})}\right) \\
        &= \frac{1}{n-1}\E\left(\sum_{i=1}^{n}\sqr{(X_i - \bar{X})}\right) \\
        &= \frac{1}{n-1}\E\left(\sum_{i=1}^{n} (X_i^2 - 2X_i\bar{X} + \bar{X}^2)\right) \\
        &= \frac{1}{n-1}\E\left(\sum_{i=1}^{n} X_i^2 - 2\bar{X}\sum_{i=1}^{n} X_i + \sum_{i=1}^{n} \bar{X}^2\right) \\
        &= \frac{1}{n-1}\E\left(\sum_{i=1}^{n} X_i^2 - 2\bar{X}(n\bar{X}) + n\bar{X}^2\right) \\
        &= \frac{1}{n-1}\E\left(\sum_{i=1}^{n} X_i^2 - n\bar{X}^2\right) \\
        &= \frac{1}{n-1}\left( \sum_{i=1}^{n} \E(X_i^2) - n\E(\bar{X}^2) \right)\\
        &= \frac{1}{n-1}\left( n(\Var(X) + (\E(X))^2) - n\left(\frac{\Var(X)}{n} + (\E(X))^2\right) \right) \\
        &= \frac{1}{n-1}\left( n\Var(X) + n(\E(X))^2 - \Var(X) - n(\E(X))^2 \right) \\
        &= \frac{1}{n-1}\left( (n-1)\Var(X) \right) \\
        &= \Var(X)
    \end{align*}
\end{proof}


\noindent Measure of Shape:
\begin{itemize}
    \item The sample skewness: 
        $$ g_1 = \frac{\frac{1}{n} \sum_{i=1}^{n}\left(y_i - \bar{y}\right)^3}{\sqr{\left[\frac{1}{n} \sum_{i=1}^{n}\sqr{\left(y - \bar{y}\right)}\right]}}$$
    \item The sample kurtosis:
        $$ g_2 = \frac{\frac{1}{n} \sum_{i=1}^{n}\left(y_i - \bar{y}\right)^4}{\sqr{\left[\frac{1}{n}\sum_{i=1}^{n}\sqr{\left(y_i - \bar{y}\right)}\right]}}$$
\end{itemize}

\noindent Interquartile range: $IQR = q(0.75) - q(0.25)$\\
\noindent Five number summary: Smallest observation, the lower quartile, the median, the upper quartile, and the largest value. $y_{(1)}, q(0.25), q(0.5), q(0.75), y_{(n)}$.

Sample Correlation:The sample correlationfor data ${(x_1, y_1), \dotsc,(x_n, y_n)}$
$$r = \frac{S_{xy}}{\sqrt{S_{xx}S_{yy}}} = \frac{\left(\sum_{i=1}^{n}\left(x_i - \bar{x}\right)\left(y_i - \bar{y}\right)\right)}{\left(\sum_{i=1}^{n} \sqr{\left({x_i - \bar{x}}\right)}\right) \left(\sum_{i=1}^{n}\sqr{\left({y_i - \bar{y}}\right)} \right)}$$


\subsection{Point Estimate and Maximum Likelihood Estimation}
A \textbf{point estimate} of a parameter is the value of a function of the observed data $y_1, \dotsc, y_n$ and other known quantities, use $\bar{\theta}$ to denote an estimate of the parameter $\theta$.
Likelihood Function for Discrete Distribution:
$$L(\theta) = L(\theta; y) = P(Y = y; \theta) \indent \text{for} \theta \in \varOmega$$
where the parameter space $\varOmega$ is the set of possible value for $\theta$.
\\\\
\noindent Maximize Likelihood Esitmation (MLE): If the value of $\theta$ which maximize the probability of observing the data $y$. Denote as $\bar{\theta}$.
Example: $Y \sim Bin(n, \theta)$
$$L(\theta) = P(Y = y; \theta) = \binom{n}{y}\theta^y(1-\theta)^{n-y} \indent \text{for} 0 \le \theta \le 1$$
Assume $y \ne 0$, $y \ne n$, let $\frac{dL(\theta)}{d\theta} = 0$.
\begin{align*}
    0 &= \binom{n}{y}\left[y\theta^{y-1}(1-\theta)^{n-y} + \theta^y(1-\theta)^{n-y-1}(-1)(n-y)\right]\\
    &= \binom{n}{y}\theta^{y - 1}(1-\theta)^{n-y-1}(y(1- \theta) - \theta(n-y))\\
&\Rightarrow y - y\theta - n\theta + y\theta = 0 \\
&\Rightarrow y = n\theta \text{\indent Sample Proportion}
\end{align*}
\\
\noindent Relative Likelihood Function: 
$$R(\theta) = \frac{L(\theta)}{L(\bar{\theta})} \indent \text{for} \theta \in \varOmega, 0 \le R(\theta) \le 1, \forall \theta \in \varOmega$$
\\
\noindent Log Likelihood Function:
$$l(\theta) = \ln L(\theta) = \log L(\theta) \indent \text{for } \theta \in \varOmega$$
Since the $\log$ function is a monotonw funciton of $L(\theta)$, two function increase and decrease toether. Both fucntions have a maximum at the same value $\theta = \hat{\theta}$
\begin{figure}[htbp]
        \centering
        \includegraphics[width = 0.4\textwidth]{images/微信图片_20260716115705_769_41.png}
\end{figure}
\\\\
Example: Likelihood Function for Poisson Distirbution
\begin{align*}
    L(\theta) &= \prod_{i=1}^{n} \frac{\theta^{y_i} e^{-\theta}}{y_i!} \\
    &= \left( \prod_{i=1}^{n} \frac{1}{y_i!} \right) \theta^{\sum y_i} e^{-n\theta} \\
    &\propto \theta^{n\bar{y}} e^{-n\theta} \quad \text{for } \theta \ge 0 \\
    \\
    l(\theta) &= \log(L(\theta)) = n\bar{y}\log\theta - n\theta + \text{constant} \\
    \frac{d l(\theta)}{d\theta} &= \frac{n\bar{y}}{\theta} - n \\
    \text{Set derivative to 0: } & \frac{n\bar{y}}{\hat{\theta}} - n = 0 \\
    &\implies \hat{\theta} = \bar{y} \quad \text{maximizes } l(\theta)
\end{align*}

\noindent Combining Likelihood Based on Independent Experiments.
Suppose$y_1$ and $y_2$ statistically independent.
$$P(Y_1 = y_1, Y_2 = y_2; \theta) = P(Y_1 = y_1; \theta)P(Y_2 = y_2;\theta)$$
Likelihood Function  $L(\theta)$ based on $y_1$ and $y_2$.
$$L(\theta) = \L_1(\theta)L_2(\theta) \indent \text{for } \theta \in \varOmega, \text{where}L_i(\theta) = P(Y_i = y_i; \theta) \text{ for } i = 1, 2,\dotsc$$
\\\\
$Y\sim Bin(n, \theta)$ and Suppose $y_1, y_2, \dotsc, y_n$ are independent random variables from Binomial distribution.
$$P(Y=y; \theta) = f(y;\theta) = \binom{n}{y}\theta^y\left(1-\theta\right)^{\theta - y} $$
Ignoring all the constant term without $\theta$.
$$L(\theta) = \theta^y (1-\theta)^(n-y), \indent \text{for } 0\le \theta \le 1, f(y_i; \theta) = \theta^{y_i}(1-\theta)^{1-y_i}$$
\begin{align*}
    L(\theta)&= \prod_{i=1}^{n}f(y_i,\theta)\\
    &= \prod_{i=1}^{n}e^{y_i}(1-\theta)^{1-y_i}\\
    &= \theta^{\sum_{i=1}^{n}y_i}(1-\theta)^{\sum_{i=1}^{n} (1- y_i)}\\
    &= \theta^y (1 - \theta)^(n-y) \indent \text{for } 0\le \theta \le 1
\end{align*}
$$l(\theta) = ln(L(\theta)) = ln\left[\theta^y(1-\theta)^{n-y}\right] = yln(\theta) + (n-y)ln(1-\theta)$$
$$\frac{dl(\theta)}{d\theta} = \frac{y}{\theta} (n-y)\frac{-1}{1-\theta} \Rightarrow y(1-\theta) - \theta(n-y) = 0 \Rightarrow y -n\theta = 0$$
Therefore $\hat{\theta} = \frac{y}{n}$



\newpage
\subsection{Likelihood Function for Continuous Distribution}
If $y_1, \dotsc, y_n$ are observed value of random sample from distribution with probability density function $f(y_i, 0)$, then
$$L(\theta) = L(\theta,y) = \prod_{i=1}^{n}f(y_i, \theta) \indent \text{for }\theta \in \varOmega$$
Example: Likelihood Function for Exponential Distribution.
Since $f(x;\theta) = \frac{1}{\theta}e^{\frac{-y_i}{\theta}}$
$$L(\theta) = \prod_{i=1}^{n}\frac{1}{\theta} e^{\frac{-y_i}{\theta}} = \frac{1}{\theta^n}\exp\left(-\frac{1}{\theta}\sum_{i=1}^{n}y_i\right) = \theta^{-n}e^{\frac{-n\bar{y}}{\theta}} \indent \text{for }\theta > 0$$
$$l(\theta) = ln(L(\theta)) = -nln(\theta) + (\frac{-n\bar{y}}{\theta}) = -n(\log\theta+\frac{\bar{y}}{\theta})$$
$$\frac{dl(\theta)}{d\theta} = \frac{-n}{\theta} + \frac{n\bar{y}}{\sqr{\theta}} = \frac{n}{\sqr{\theta}}\left(\bar{y} - \theta\right) = 0 \implies \hat{\theta} = \bar{y} \text{ is the maximum likelihood estimate of } \theta$$\\\\
\noindent Example Likelihood Funciton for Gaussian Distribution.
\begin{align*}
    L(\theta) &= L(\mu, \sigma) = \prod_{i=1}^{n}f(y_i; \mu, \sigma) = \prod_{i=1}^{n}\frac{1}{\sqrt{2\pi}\sigma}\exp\left[\frac{1}{-2\sqr{\sigma}}\sqrt{(y_i - \mu)}\right]\\
&= (2\pi)^{\frac{-\pi}{2}}\sigma^{-n}\exp\left[\sum_{i=1}^{n}\frac{-1}{2\sqr{\sigma}}\sqr{(y_i - \mu)}\right]
\end{align*}
Ignoring all the constant term without $\mu$ and $\theta$.
$$L(\theta) = L(\mu,\sigma) = \sigma^{-n}\exp\left[\frac{-1}{2\sqr{\sigma}}\sum_{i=1}^{n}\sqr{(y_i - \mu)}\right]$$
Since
$$\sum_{i=1}^{n}(y_i - \bar{y}) = \sum_{i=1}^{n}(y_i) - \sum_{i=1}^{n}(\bar{y}) = \sum_{i=1}^{n}(y_i) - ny_i = 0$$
and
\[
\sum_{i=1}^{n}\sqr{(y_i - \mu)} = \sum_{i=1}^{n}\sqr{(y_i - \bar{y} + \bar{y} - \mu)\Rightarrow \sum_{i=1}^{n}\sqr{(y_i - \bar{y})} - 2\sum_{i=1}^{n}(y_i - \bar{y})(\bar{y} - \mu)} - \sum_{i=1}^{n}\sqr{(\bar{y} - \mu)} \Rightarrow \sum_{i=1}^{n}\sqr{(y_i - \bar{y})} - \sum_{i=1}^{n}\sqr{(\bar{y} - \mu)} = 0
\]
\[
L(\theta) = \sigma^{-n}\exp\left[\frac{-1}{2\sqr{\sigma}}\sum_{i=1}^{n}\sqr{(y_i - \bar{y})}\right]\exp\left[\frac{-1}{2\sqr{\sigma}}\sum_{i=1}^{n}\sqr{(\bar{y} - \mu)}\right]
\]
The log likelihood funciton is
\[
l(\theta) = l(\mu, \sigma) = -n\log\sigma - \frac{1}{2\sqr{\sigma}}\sum_{i=1}^{n}\sqr{(y_i - \bar{y})} - \frac{1}{2\sqr{\sigma}}n\sqr{(\bar{y} - \mu)} \indent \text{for } \mu \in \R \text{ and for } \sigma > 0
\]
Taks partial derivative to each variables
\[
\frac{\partial(\theta)}{\partial \mu} = \frac{-1}{2\sqr{\sigma}}2n(\bar{y} - \mu) = \frac{-n}{\sqr{\sigma}}(\bar{y} - \mu) = 0 \Rightarrow \hat{\mu} = \bar{y}
\]
Substitute $\hat{\mu} = \bar{y}$ into partial equation and complete expansion of $\sqr{(y_i - \mu)}$

\begin{align*}
    \frac{\partial l(\theta)}{\partial \sigma} &= \frac{-n}{\sigma} + \frac{1}{\sigma ^3}\left(\sum_{i=1}^{n}\sqr{(y_i - \bar{y})} - 2(y_i - \bar{y})(\bar{y} - \mu) + \sqr{(\bar{y} - \mu)}\right) \\
    &= \frac{-n}{\sigma} + \frac{1}{\sigma^3}(\sum_{i=1}^{n}\sqr{(y_i - \bar{y})})\\
\end{align*}
\[
\frac{\partial l(\theta)}{\partial \sigma} = \frac{-n}{\sigma} + \frac{1}{\sigma^3}(\sum_{i=1}^{n}\sqr{(y_i - \bar{y})}) = 0
\]
$$\frac{n}{\sigma} = \frac{1}{\sigma^3}\left(\sum_{i=1}^{n}\sqr{(y_i - \bar{y})}\right) \Rightarrow \frac{n}{\sigma} \sigma^3 = \sum_{i=1}^{n}\sqr{(y_i -\bar{y})} \Rightarrow \sqr{\sigma} = \frac{\sum_{i=1}^{n}\sqr{(y_i - \bar{y})}}{n}\Rightarrow \hat{\sigma} = \sqrt{\frac{\sum_{i=1}^{n}\sqr{(y_i - y\bar{y})}}{n}}$$

\subsection{Likelihood Functions for Multinomial Models}
Since the multinomial joint probability function is
\[
f(y_1, y_2, \dotsc, y_k;\theta) = \frac{n!}{y_1!y_2!\dotsc y_k!}\prod_{i=1}^{k}\theta_i^{y_i} \indent \text{for } y_i = 0, 1, \dotsc \text{where }\sum_{i=1}^{k}y_i = n
\]
So the likelihood function for $\theta = (\theta_1, \dotsc, \theta_k)$ based on data $y_1, \dotsc, y_k$ is given by
\[
L(\theta) = L(\theta_1, \dotsc, \theta_k) = \frac{n!}{y_1! \dotsc y_k!}\theta_i^{y_i}
\]

\noindent Ignoring all the constant term withou $\theta$.
\[
L(\theta) = \prod_{i=1}^{k}\theta_i^{y_i}
\]
The log likelihood function is
\[
l(\theta) =\sum_{i=1}^{k}y_i \log\theta_i
\]
Therefore the maximum likelihood estimation of $\theta_1, \dotsc, \theta_k$ are
\[
\hat{\theta} = \frac{y_i}{n}
\]

\noindent Example: Suppose 
\[
P(Y_1 = y_1, Y_2 = y_2, Y_3 = y_3;\theta) =\frac{n!}{y_1!\dotsc y_3!}y_1^{\theta_1}y_2^{\theta_2}y_3^{\theta_3}
\]
For $\theta_1  = \sqr{\alpha}, \theta_2 = 2\alpha(1-\alpha), \theta_3 = \sqr{1-\alpha}$ where $0\le \alpha \le 1$.
Then the likelihood function is
\[
L(\alpha) = \frac{n!}{y_1!y_2!y_3!}(\sqr{\alpha})^{y_1}(2\alpha(1-\alpha))^{y_2}\left[\sqr{1-\alpha}\right]^{y_3}
\]
\[
L(\alpha) = = 2^{y_2}a^(2y_1 + y_2)(1-\alpha)^(y_2 + 2y_3)
\]
The log likelihood function is
\[
l(\alpha) = (2y_1 + y_2)\log\alpha + (y_2 + 2y_3)\log(1-\alpha) \Rightarrow \frac{dl(\alpha)}{d\alpha} = \frac{2y_1 + y_2}{\alpha} + \frac{y_2 + 2y_3}{1-\alpha} =0
\] 
Then 
\[
\alpha = \frac{2y_1 + y_2}{2y_1 + y_2 + y_3} = \frac{2y_1 + y_2}{2n}
\]
is the maximum likelihood of $\alpha$.
\newpage
\noindent Theorem: If $\bar{\theta} = \left(\bar{\theta}_1, \dotsc, \bar{\theta}_k\right)$ is the maximum likelihood estimate of $\theta = \left(\theta_1, \dotsc, \theta_k\right)$, then $g(\bar{\theta})$ is the maximum likelihood estimate of $g(\theta)$.

\section{Estimation}
\subsection{Esitmator and Sampling Distribution}
A estimator $\tilde{\theta}$ is a random variabel which is a function $\tilde{\theta} = g(Y_1, \dotsc, Y_n)$ fo the random variables $Y_1, \dotsc, Y_n$. The distribution of $\theta$ is called the sampling distirbution of the estimator.\\

\noindent Exmaple:Suppose to take a random sample $Y_1, \dotsc, Y_n$ to estimate the unknown mean $\mu$ where $Y_i \sim G(\mu, \sigma), i = 1, \dotsc, n$ and the maximum likelihood estimation of $\mu$ is 
\[
\tilde{\mu} = \bar{Y} = \frac{1}{n}\sum_{i=1}^{n}Y_i
\]
\[
\Var(\bar{Y}) = \Var(\frac{1}{n}\sum_{i=1}^{n}Y_i) = \frac{1}{\sqr{n}}\sum_{i=1}^{n}\Var(Y_i) = \frac{1}{\sqr{n}}(n\sqr{\sigma}) = \frac{\sqr{\sigma}}{n} \Rightarrow sd(\bar{Y}) = \frac{\sigma}{\sqrt{n}}
\]
then $\tilde{\mu} = \bar{Y}\sim G(\mu, \frac{\sigma}{\sqrt{n}})$.Dtermine how often of estimator $\tilde{\mu} = \bar{Y}$ within 0.01.
\[
P(\abs{\tilde{\mu} - \mu} \le 0.01) = P(\mu - 0.01 \le \bar{Y} \le \mu + 0.01) = P\left(\frac{-0.01}{\frac{\sigma}{\sqrt{n}}} \le \frac{\bar{Y} - \mu}{\frac{\sigma}{\sqrt{n}}} \le \frac{0.01}{\frac{\sigma}{\sqrt{n}}}\right) \Rightarrow P\left(\mu - 0,01 \le z \le \mu + 0.01\right) \text{where } z\sim G(0, 1) 
\]
Suppose $\sigma = 0.07$ and $n = 50$, then
\[
P(\abs{\tilde{\mu} - \mu}\le 0.01) = P(\frac{-0.01\sqrt{50}}{0.07} \le z \le \frac{0.01\sqrt{50}}{0.07})  = P(-1.01\le z \le 1.01) = 0.688
\]
\subsection{Confidence Intervals amd Pivotal Quantities}
Example: Gaussian Distribution with unknown mean and known standard deviation
Suppose that the observation$y_1, \dotsc, y_n$ are random sample from $G(\mu, 2)$ distribution and $\mu = \E(Y_i)$ is unknown but $sd(Y_i) = 2$The maximum likelihood of $\mu$ is $\hat{\mu} = \bar{y}$ and the MLE of $\mu$
\[
\left[\bar{y} - 1.96(2)/\sqrt{n}, \bar{y} + 1.96(2/\sqrt{n})\right]
\] 
The interval is random interval which both endpoints are random variables. Since $Y \sim G(\mu, \frac{2}{\sqrt{n}})$
\begin{align*}
    P\left(\mu \in \left[\bar{Y} - 1.96(2)/\sqrt{n}, \bar{Y}+ 1.96(2)/\sqrt{n}\right]\right) &= P(\bar{Y} - 1.96(2)/\sqrt{n} \le \mu \bar{Y} + 1.96(2)/\sqrt{n})\\
    &= P(-1.96\le \frac{\bar{y} -\mu}{2/\sqrt{n}} \le 1.96)\\
    &= P(-1.96 \le Z \le 1.96) \text{where }Z\in G(0,1)\\
    &= 0.95
\end{align*}
\noindent Interval Estimator: The interval estimator $\left[L(Y),U(Y)\right]$ has property that
\[
P(\theta \in \left[L(Y), U(Y)\right]) = P(L(Y) \le \theta U(Y)) = p
\]
The interval estimate $\left[L(y), U(y)\right] a100p\%$ confidence interval for $\theta$ and $p$ is condifence coefficient\\\\
\noindent Pivotal Quantity: $Q = Q(Y;\theta)$ is a funciton of data $Y$ and the unknown parameter $\theta$ such that the distirbution of the random variable $Q$ is fully known. The probability statements such as $P(Q\le b)$ and $P(Q\le a)$depend on a and b but not on $\theta$ or any unknown information.

\subsubsection{Asymptotic Gaussian Pivotal Quantities}
Suppose $\bar{\theta}$ is point estimator of the unknown parameter of $\theta$. Supppose that the Central Limit Theorem can be used to obtain the result that
\[
\frac{\bar{\theta} - \theta}{g(\theta) - \sqrt{n}}
\]
has approximately $G(0,1)$ distribution for large n where $\E(\bar{\theta}) = \theta$ and $sd(\bar{\theta}) = h(\theta)/\sqrt{n}$ for some real valued $g(\theta)$.We replace $\theta$ by $\bar{theta}$ in the denominator then
\[
Q_n(\bar{\theta},\theta) = \frac{\bar{\theta} - \theta}{g(\bar{\theta})/\sqrt{n}}
\]
alse has approximately a $G(0,1)$ distribution for large n. Therefore $Q(\bar{\theta},\theta)$ is an asymptotic Gaussian pivotal quantity which can be used to construct approximate confidence interval for $\theta$.

\subsection{The chi-squared and t Distributions}
\subsubsection{Gamma Function}

\[
\Gamma(\alpha) = \uint{0}{\infty}{y^{\alpha - 1}e^{-y}}{y}
\]
\begin{itemize}
    \item $\Gamma(\alpha) = (\alpha - 1)\Gamma(\alpha - 1)$
    \item $\Gamma(\alpha) = (\alpha - 1)! , \text{for } \alpha \in \N, a > 0$
    \item $\Gamma(\frac{1}{2}) = \sqrt{n}$
\end{itemize}
The $\sqr{\chi}(k)$ distribution is a continuous family of distribution $(0, \infty)$ with probability density funciton of the form
\[
f(x;k) = \frac{1}{2^{\frac{k}{2}}\Gamma(\frac{k}{2})}x^{\frac{k}{2} - 1}e^{\frac{-x}{2}} \indent \text{for } x > 0, \text{ where } k \in \left\{k = 1, 2, \dotsc\right\}
\]
Denote $X \sim \sqr{\chi(k)}$
\begin{itemize}
    \item for $k = 2$ substitute into equation $f(x; 2) = \frac{1}{2\Gamma(1)}x^0e^{\frac{-x}{2}} = \frac{1}{2}e^{\frac{-x}{2}}$ equivalent to $Exponential(2)$ 
    \item for $k > 2$ then $pdf$ is unimodal with maximum value at $x= k - 2$.
    \item for $k \ge 30$  then $pdf$ resembles that of $N(k, 2k) pdf$ 
\end{itemize}
If $X \sim \sqr{\chi}$, Then
\[
\E(X) = k  \text{ and } \Var(X) = 2k
\]
\begin{align*}
    \E(X^j) &= \uint{0}{\infty}{x^j\frac{1}{2^{\frac{k}{2}\Gamma(\frac{k}{2})}}x^{(\frac{k}{2} -1)}e^{\frac{k}{2}}}{x}\\
    &= \frac{1}{\Gamma(\frac{k}{2})}\uint{0}{\infty}{\frac{1}{2^{\frac{k}{2}}}x^{(\frac{k}{2}+j - 1)}e^{\frac{k}{2}}}{x} \indent \text{Let } y = \frac{x}{2} \Rightarrow dx = 2dy\\
    &= \frac{1}{\Gamma(\frac{k}{2})}\uint{0}{\infty}{2^{-\frac{k}{2}}(2y)^{\frac{k}{2} + j - 1}2}{y}\\
    &= \frac{2^j}{\Gamma(\frac{k}{2})}\uint{0}{\infty}{y^(\frac{k}{2}+ j - 1)}{y}\\
    &= \frac{2^j \Gamma(\frac{k}{2} + j)}{\Gamma(\frac{k}{2})}
\end{align*}
For $j = 1$, then 
\[
\E(X) = \frac{2\Gamma(\frac{k}{2} +1)}{\Gamma{\frac{k}{2}}} = \frac{2 \frac{k}{2} \Gamma(\frac{k}{2})}{\Gamma(\frac{k}{2})} = k
\]
\[
\E(\sqr{X}) = \frac{2^2 \Gamma(\frac{k}{2} + 2)}{\Gamma(\frac{k}{2})} = \frac{4 \Gamma((\frac{k}{2} + 1) + 1)}{\Gamma(\frac{k}{2})} = 4(\frac{k}{2} + 1)(\frac{k}{2}) = k(k+ 1)
\]

\begin{align*}
    \uint{0}{\infty}{\frac{1}{2^{\frac{k}{2}}\Gamma(\frac{k}{2})}x^{\frac{k}{2}-1}e^{\frac{-x}{2}}}{x} &= \uint{0}{\infty}{\frac{1}{2^{\frac{k}{2}}\Gamma(\frac{k}{2})} (2y)^{\frac{k}{2} -1}e^{-y}2y}{y}\\
    &= \frac{2^{\frac{k}{2} - 1 + 1}}{2^{\frac{k}{2}}\Gamma(\frac{k}{2})}\uint{0}{\infty}{y^{\frac{k}{2} -1} e^{\frac{-y}{2}}}{y}\\
    &= \frac{\Gamma(\frac{k}{2})}{\Gamma(\frac{k}{2})}\\
    &= 1
\end{align*}
mgf of $\sqr{\chi}$ distirbution:
\begin{align*}
    M(t) = \E(e^{tX}) &= \uint{0}{\infty}{e^{tx}f(x, k)}{x} \\
&= \uint{0}{\infty}{e^{tx}\frac{1}{2^{\frac{k}{2}}\Gamma(\frac{k}{2})}x^{\left(\frac{k}{2} - 1\right)}e^{\frac{-x}{2}}}{x}\\
&= \uint{0}{\infty}{\frac{1}{2^{\frac{k}{2}} \Gamma(\frac{k}{2})} x^{\left(\frac{k}{2} - 1\right)} e^{-x\left(1- 2t\right)}}{x} \indent \text{let }u = x\frac{1-2t}{2} \text{ and } \frac{2u}{1-2t}\Rightarrow dx = \frac{2du}{1-2t}\\
&= \uint{0}{\infty}{\frac{1}{2^{\frac{k}{2}}\Gamma(\frac{k}{2})}\left(\frac{2u}{1-2t}\right)^{\frac{k}{2} -1 } e^{-u}}{u}\\
&= \uint{0}{\infty}{\frac{1}{2^{\frac{k}{2}}\Gamma(\frac{k}{2})}\left(\frac{2}{1-2t}\right)^{\frac{k}{2}}u^{\frac{k}{2} - 1}e^{-u}}{u}\\
&= \frac{1}{\Gamma(\frac{k}{2})(1-2t)^\frac{k}{2}}\uint{0}{\infty}{u^{\frac{k}{2} - 1}e^{-u}}{u}\\
&= \frac{\Gamma(\frac{k}{2})}{\Gamma(\frac{k}{2})(1-2t)^{\frac{k}{2}}}\\
&= (1 - 2t)^{-\frac{k}{2}} \indent \text{for } t < \frac{1}{2}
\end{align*}
\noindent Theorem: The additive of $\sqr{\chi}$ variables.\\
Let $W_1, \dotsc, W_n$ be independent random variables with $W_i\sim \sqr{\chi}(k_i)$. Then
\[
S = \sum_{i=1}^{n}W_i \sim \sqr{\chi}\left(\sum_{i=1}^{n}k_i\right)
\]
\begin{proof}
    Let $S = \sum_{i=1}^{n}W_i$. Since$W_1, \dotsc, W_n$ are independent random variables  with $W_i \sim \sqr{\chi}(k_i)$, then the $mgf$ of $S$
    \begin{align*}
        M_S(t) &= \prod_{i=1}^{n}M_{W_i}(t)\\
        &= \prod_{i=1}^{n}(1-2t)^{\frac{-k_i}{2}}\\
        &= (1-2t)^{\left(\sum_{i=1}^{n}k_i\right)}\\
        &= (1-2t)^{\frac{-1}{2}\sum_{i=1}^{n}k_i}
    \end{align*}
Thus $S = \sum_{i=1}^{n}W_i \sim \sqr{\chi}(k_i)$
\end{proof}

\noindent Theorem: If $Z \sim G(0,1)$ then the distribution of $W=\sqr{Z}$ is $\sqr{\chi}(1)$
\begin{proof}
    Suppose that $W = \sqr{Z}$ where $Z\sim G(0,1)$. Let $\Phi$ represent the $cdf$ of $G(0,1)$random variables and let $\phi$ represent the $pdf$ of $G(0,1)$ of random variable.
    For $w>0$
\[
P(W\ge w) = P(-\sqrt{w}\le z \le \sqrt{w}) = \Phi(\sqrt{w}) - \Phi(-\sqrt{w})
\]
Then the $pdf$ of $W$ is

\begin{align*}
    \frac{d}{dw}P(W\le w) &= \frac{d}{dw}\left[\Phi(\sqrt{w}) - \Phi(-\sqrt{w})\right] \\
&= \left[\phi(\sqrt{w}) - \phi(-\sqrt{w})\right]\left(\frac{1}{2}w^{\frac{-1}{2}}\right)\\
&= 2\left[\frac{1}{\sqrt{2\pi}}e^{\frac{-(\sqr{\sqrt{w}})}{2}}\frac{1}{2\sqrt{w}}\right]\\
&= \left[\frac{1}{\sqrt{2\pi}}e^{\frac{-w}{2}}w^{\frac{-1}{2}}\right]\\
&= \frac{1}{2^{\frac{1}{2}}\Gamma(\frac{1}{2})}w^{\frac{1}{2} - 1}e^{\frac{-w}{2}}
\end{align*}
Thus the probability density function of a $\sqr{\chi}(1)$ random variable as required\\
\end{proof}
\noindent Corollary: If $Z_1, \dotsc, Z_n$ are mutually independent $G(0,1)$ random variables and $S = \sum_{i=1}^{n}\sqr{Z_i}$, then $S\sim \sqr{\chi}(n)$
\begin{proof}
    Since $Z_i \sim G(0,1)$ then by previous theorem $W = \sqr{Z_i} \sim \sqr{\chi}(1)$. According to the additive of $\sqr{\chi}$ variables, $W_1, \dotsc, W_n$ are mutually independent random variable, let $S = \sum_{i=1}^{n}W_i$. Therefore
    $$S = \sum_{i=1}^{n}W_i \sim \sqr{\chi}\left(\sum_{i=1}^{n}1\right) = \sqr{\chi}(n)$$ 
\end{proof}

\noindent Theroem: If $X \sim \sqr{\chi}(n_1), Y\sim \sqr{\chi}(n_2)$ also X and Y are mutually independent. Then
\[
X + Y \sim \sqr{\chi}(n_1 + n_2)
\]
\subsubsection{Student's $t$ Distribution}
The $pdf$ of Student's $t$ distribution:
\[
f(t;k) = c_k\left(1+\frac{\sqr{t}}{k}\right)^{-\frac{k+1}{2}} \indent \text{for } t\in \R \text{ and } k = 1, 2,\dotsc
\]
where
\[
c_k= \frac{\Gamma(\frac{k+1}{2})}{\sqrt{k\pi}\Gamma(\frac{k}{2})}
\]
The parameter $k$ is the degrees of freedom. Denote as $T\sim t(k)$ to indicate that the random variable $T$ has $t$ distribution with $k$ degrees of freedom.For smaller k(such as $k=2$), 
the peak of $t$ distribution is lower and has "fatter tails" at both ends.
Implies that the probability of extreme value is higher than standard normal distribution.
as $t -> \infty, f(t, k) -> \frac{1}{\sqrt{2\pi}}e^{-\frac{\sqr{x}}{2}}$

\begin{figure}[htbp]
    \centering
    \includegraphics[width = 0.4\textwidth]{images/微信图片_20260721131020_1_17.png}
\end{figure}
\newpage
\noindent Theroem:Suppose $Z \sim G(0,1)$ and $U\sim \sqr{\chi}(k)$ independently. Let
\[
T = \frac{Z}{\sqrt{\frac{U}{k}}}
\]
Then T has a Student's t distribution with k degrees of freedom.

\subsection{Likelihood-Based Confidence Intervals}
The relative likelihood
\[
R(\theta) = \frac{L(\theta)}{L(\hat{\theta})}\indent \text{for }\theta \in \varOmega
\]
is the function of the maximum lieklihood estimate $\hat{\theta}$. Use random variable $\bar{\theta}$ to substitute $\hat{\theta}$ and define $\lambda(\theta)$
\[
\lambda(\theta) = -2\log\left[\frac{L(\theta)}{L(\bar{\theta})}\right]
\]
where $\bar{\theta}$ is the maximum likelihood estimator and $\lambda(\theta)$ is the likelihood ratioi statistic.\\\\
Theorem: If $L(\theta)$ is based on $Y = (Y_1, \dotsc, Y_n)$ a ranodm sample of size $n$ and if $\theta$ is the true value of the each parameter, then thedsitribution of $\lambda(\theta)$ converges to $\sqr{\chi}(1)$ distribution as $n -> \infty$.
\\

\noindent Theorem: A $100p\%$ lieklihood interval is an approximate $100q\%$ confidence interval where $q= wP(Z \ge \sqrt{-2 log p} - 1)$ and $Z \sim N(0,1)$.
\begin{proof}
    A $100p\%$ likelihood interval is defined by $\{\theta \mid R(\theta) \ge p\}$, which can be rewritten as:
    \[
    \{\theta \mid R(\theta) \ge p\} = \left\{\theta \mid -2\log\left[\frac{L(\theta)}{L(\hat{\theta})}\right] \le -2\log p\right\}
    \]
    \begin{align*}
    P(\Lambda(\theta) \le -2\log p) &= P\left(-2\log\left[\frac{L(\theta)}{L(\hat{\theta})}\right] \le -2 \log p\right) \\
    &= P(W \le -2\log p) \indent \text{for } w\sim \sqr{\chi}(1) \text{ from previous theorem}\\
    &= P(-\sqrt{-2\log p} \le Z \le \sqrt{-2 \log p}) \indent \text{since } Z \sim N(0,1)\\
    &= P(Z \le \sqrt{-2 \log p}) - \left[1 - P(- \sqrt{-2 \log p} \le Z)\right]\\
    &= 2P(Z \le\sqrt{-2 \log p}) - 1
\end{align*}
\end{proof}
Example: If $p= 0.1$, then 
\begin{align*}
    q &= 2P\left(\sqrt{-2 \log (0.1)}\right) - 1\indent \text{where }Z \sim N(0,1)\\
    &= 2P(Z\le 2.15) - 1\\
    &= 0.96844
\end{align*}

\subsection{Condifence Intervals for Parameters in the $G(\mu, \sigma) Model$}
\end{document}